% SC3020 --- Chapter 1: Introduction to Data Analytics and Mining (0->100 ground-up)
\chapter{Introduction to Data Analytics and Mining}
\label{ch:intro}

\begin{quote}
\emph{``Data is the new oil, but analytics is the refinery.''} Mining turns raw records into decisions. This chapter builds the why and what of the field from everyday examples to a precise process model.
\end{quote}

\noindent\fbox{\parbox{\dimexpr\linewidth-2\fboxsep-2\fboxrule}{%
\textbf{From zero: no analytics background assumed.} We start from a shopkeeper's ledger, define \emph{data, information, knowledge}, introduce types of data you will meet in every project, then formalise the CRISP-DM lifecycle. Pictures precede formalism.}}
\vspace{0.4em}

\section*{Learning Outcomes}
After this chapter you will be able to:
\begin{enumerate}[label=(L\arabic*)]
    \item distinguish data, information, and knowledge and explain the analytics pipeline;
    \item classify data types (nominal, ordinal, interval, ratio; structured vs.\ unstructured);
    \item describe the six phases of CRISP-DM and their iterative nature;
    \item identify ethical issues (privacy, bias, fairness) and relevant regulations;
    \item scope a small analytics project using a CRISP-DM checklist.
\end{enumerate}

% ============================================================
\section{Why Analytics? From Ledger to Decisions}
\label{sec:why-analytics}
\paragraph{Intuition.} A neighbourhood bakery writes daily sales in a notebook: time, item, price, customer. The notebook is \emph{data} (raw facts). Counting ``croissants sell twice as fast on Saturdays'' is \emph{information} (processed data). Deciding ``bake 30 extra croissants Saturday morning'' is \emph{knowledge} (actionable information). \emph{Analytics} is the refinery between them; \emph{data mining} is the set of algorithms that find the non-obvious patterns.

\paragraph{Formal definitions.} Let a \emph{dataset} be a collection $\mathcal{D}=\{x_i\}_{i=1}^n$ where each record $x_i\in\mathcal{X}$ describes an entity via attributes (features). An \emph{analytics pipeline} transforms $\mathcal{D}$ into a model $M$ and then into a decision $d$ that maximises expected utility under a business objective. Two classic frameworks organise this: \emph{KDD} (Knowledge Discovery in Databases, Fayyad et al.) and \emph{CRISP-DM} (Cross-Industry Standard Process for Data Mining). KDD is academic (selection $\to$ preprocessing $\to$ transformation $\to$ mining $\to$ interpretation); CRISP-DM is industrial and we adopt it here.

% ============================================================
\section{Types of Data}
\label{sec:data-types}
\paragraph{Intuition.} Not all columns behave alike. You can average a temperature but not a postcode, and you can rank ``small $<$ medium $<$ large'' but not compute ``medium minus small.'' Using the wrong operation is a silent bug.

\paragraph{Formal taxonomy.} Stevens' scales:
\begin{description}[leftmargin=1.2em, itemsep=0.3em]
    \item[Nominal] Categories without order (e.g., blood type, country). Only $=$, $\neq$, mode.
    \item[Ordinal] Ordered categories without meaningful gaps (e.g., rating $1$--$5$, education level). Adds $<$ and median.
    \item[Interval] Numeric with arbitrary zero (e.g., Celsius, calendar year). Differences meaningful, ratios not ($20^\circ$ is not ``twice as hot'' as $10^\circ$).
    \item[Ratio] Numeric with absolute zero (e.g., age, price, Kelvin). All arithmetic valid.
\end{description}
Orthogonally, \emph{structured} data sits in tables (rows $=$ records, columns $=$ attributes); \emph{semi-structured} has self-describing tags (JSON, XML); \emph{unstructured} has no fixed schema (text, images, audio). Most mining value now lies in the latter two, but almost always begins by \emph{structuring} them (e.g., bag-of-words).

\begin{figure}[h]\centering
\begin{tikzpicture}[
  st/.style={draw, thick, rounded corners=3pt, align=center, font=\scriptsize, minimum width=2.6cm, minimum height=0.9cm},
  arrow/.style={-{Latex[length=2.5mm]}, thick, color=blue!60!black}
]
\node[st, fill=blue!14] (nom) at (0,0) {Nominal\\{\tiny blood type, colour}};
\node[st, fill=orange!16] (ord) at (3.6,0) {Ordinal\\{\tiny rating, grade}};
\node[st, fill=green!14] (int) at (7.2,0) {Interval\\{\tiny $^\circ$C, year}};
\node[st, fill=violet!14] (rat) at (10.8,0) {Ratio\\{\tiny age, price, Kelvin}};
\draw[arrow] (nom)--(ord); \draw[arrow] (ord)--(int); \draw[arrow] (int)--(rat);
\node[font=\tiny, fill=yellow!12, rounded corners=2pt, inner sep=2pt] at (5.4,-0.95) {increasing information $\to$ more operations allowed (mode $\to$ median $\to$ mean $\to$ ratio)};
\draw[decorate, decoration={brace, amplitude=4pt}, thick, gray!70] (0,-1.35)--(10.8,-1.35) node[midway, below=4pt, font=\tiny] {Stevens' scales (nominal weakest, ratio strongest)};
\node[st, fill=teal!12, minimum width=3.8cm] (str) at (2.4,-2.4) {Structured (tables)};
\node[st, fill=red!12, minimum width=3.8cm] (semi) at (6.6,-2.4) {Semi-structured (JSON)};
\node[st, fill=yellow!16, minimum width=3.8cm] (un) at (10.8,-2.4) {Unstructured (text/image)};
\end{tikzpicture}
\caption{Data type hierarchy. Higher scales inherit operations of lower ones. Structured vs.\ semi/unstructured is an orthogonal axis about schema regularity.}
\label{fig:data-types}
\end{figure}

% ============================================================
\section{CRISP-DM: The Mining Lifecycle}
\label{sec:crispdm}
\paragraph{Intuition.} Without a process, projects wander: a brilliant model that solves the wrong question, or a perfect question with no deployable answer. CRISP-DM prevents that by forcing you to define success \emph{before} touching data.

\paragraph{Six phases.} (1) \textbf{Business Understanding}: translate a business goal (``reduce churn 10\%'') into a data-mining goal (``predict P(churn $|$ behaviour)''). Deliverable: success criteria. (2) \textbf{Data Understanding}: collect, describe, explore, verify quality. (3) \textbf{Data Preparation} (typically 60--80\% of effort): cleaning, integration, transformation, reduction (Ch.~2). (4) \textbf{Modeling}: select and calibrate algorithms (Ch.~3--6). (5) \textbf{Evaluation}: test against business criteria, not just accuracy (Ch.~7). (6) \textbf{Deployment}: integrate, monitor, maintain. Crucially the process is \emph{cyclic}: failure at Evaluation loops back to Business Understanding; insights in Modeling reveal new data needs.

\begin{figure}[h]\centering
\begin{tikzpicture}[every node/.style={font=\scriptsize}, scale=0.95]
\def\r{2.6}
\foreach \a/\col/\txt in {90/blue!16/Business\\Understanding,30/green!16/Data\\Understanding,330/orange!16/Data\\Preparation,270/violet!14/Modeling,210/red!12/Evaluation,150/teal!14/Deployment} {
  \node[draw, thick, fill=\col, rounded corners=4pt, align=center, minimum width=2.1cm, minimum height=1.0cm] at (\a:\r) {\txt};
}
\draw[-{Latex[length=3mm]}, thick, blue!60!black] (90:\r) ++(18:1.05) arc (18:-42:2.95);
\draw[-{Latex[length=3mm]}, thick, blue!60!black] (30:\r) ++(-42:1.05) arc (-42:-102:2.95);
\draw[-{Latex[length=3mm]}, thick, blue!60!black] (330:\r) ++(-102:1.05) arc (-102:-162:2.95);
\draw[-{Latex[length=3mm]}, thick, blue!60!black] (270:\r) ++(-162:1.05) arc (-162:-222:2.95);
\draw[-{Latex[length=3mm]}, thick, blue!60!black] (210:\r) ++(-222:1.05) arc (-222:-282:2.95);
\draw[-{Latex[length=3mm]}, thick, blue!60!black] (150:\r) ++(-282:1.05) arc (-282:-342:2.95);
% inner feedback arrows
\draw[-{Latex[length=2mm]}, thick, dashed, red!60!black] (1.2,-1.6) .. controls (0,-0.6) .. (-0.8,1.2) node[midway, right, font=\tiny, fill=white] {iterate};
\node[font=\scriptsize\bfseries, fill=yellow!12, rounded corners=2pt, inner sep=3pt] at (0,0) {CRISP-DM cycle};
\node[font=\tiny, fill=gray!8, rounded corners=2pt, inner sep=2pt] at (0,-3.7) {Outer ring = forward flow; dashed = common feedback loops. Data Preparation dominates effort.};
\end{tikzpicture}
\caption{CRISP-DM cycle. The outer ring is the nominal sequence; inner feedback edges are the norm in practice, not the exception.}
\label{fig:crispdm}
\end{figure}

\begin{lstlisting}[language=Python, caption={Classifying columns before modelling avoids illegal operations.}, label={lst:types}]
import pandas as pd
df = pd.DataFrame({"colour":["red","blue","red"], "rating":[1,3,5],
                   "temp_C":[10,20,30], "price":[3.5,4.0,5.5]})

# Stevens checks: what is meaningful?
print(df["colour"].mode())          # OK for nominal
print(df["rating"].median())        # OK for ordinal+
print((df["temp_C"].mean(), df["price"].mean()))  # both numeric
# Forbidden: ratio on interval
# print("20C twice 10C?", 20/10)   # meaningless for Celsius

# Schema check
print(df.dtypes)                    # structured table already
# For JSON column: pd.json_normalize(df["json_col"])
\end{lstlisting}

% ============================================================
\section{Ethics, Privacy, and Governance}
\label{sec:ethics}
\paragraph{Intuition.} The bakery knows who buys what and when. Publishing ``customer X buys cake every Friday'' may be harmless; linking it to health records is not. All analytics trades off insight against intrusion.

\paragraph{Framework.} Three pillars: \emph{privacy} (control over personal data --- GDPR, PDPA Singapore, consent, anonymisation, $k$-anonymity), \emph{bias and fairness} (data reflects history; models amplify it --- e.g., hiring data biased against a group yields discriminatory predictions; mitigate via balanced sampling, fairness metrics, audit), and \emph{accountability} (provenance, reproducibility, explainability, human-in-the-loop for high-stakes decisions). A useful test before deployment: would the affected person, told how the model uses their data, consider it fair? If not, revisit Business Understanding.

\begin{figure}[h]\centering
\begin{tikzpicture}[every node/.style={font=\scriptsize}]
\node[draw, thick, fill=blue!12, rounded corners=4pt, minimum width=3.2cm, minimum height=1.1cm, align=center] (a) at (0,0) {Privacy\\{\tiny GDPR, PDPA, $k$-anonymity}};
\node[draw, thick, fill=orange!16, rounded corners=4pt, minimum width=3.2cm, minimum height=1.1cm, align=center] (b) at (4.5,0) {Fairness / Bias\\{\tiny representation, measurement}};
\node[draw, thick, fill=green!14, rounded corners=4pt, minimum width=3.2cm, minimum height=1.1cm, align=center] (c) at (9.0,0) {Accountability\\{\tiny provenance, explainability}};
\draw[{Latex[length=2.5mm]}-{Latex[length=2.5mm]}, thick, gray!60] (a)--(b);
\draw[{Latex[length=2.5mm]}-{Latex[length=2.5mm]}, thick, gray!60] (b)--(c);
\draw[thick, dashed, violet!60!black] (c) .. controls (4.5,-1.6) .. (a);
\node[font=\tiny, fill=yellow!12, rounded corners=2pt, inner sep=2pt] at (4.5,-1.15) {ethical review at every CRISP-DM phase, not only at deployment};
\end{tikzpicture}
\caption{Ethics triangle. The three concerns interact; a fix for one can harm another (e.g., removing a sensitive attribute may hide bias rather than remove it).}
\label{fig:ethics}
\end{figure}

% ============================================================
% ============================================================
\section{Analytics versus Mining versus Science}
\label{sec:vs-science}
\paragraph{Intuition.} Vendors conflate terms. Think of a hierarchy: \emph{data analytics} is the umbrella (any analysis of data for insight), \emph{data mining} is the algorithmic core (pattern discovery at scale, often without a pre-stated hypothesis), and \emph{data science} adds engineering and communication (pipelines, deployment, storytelling). A data analyst asks ``what happened?'', a miner asks ``what pattern holds generally?'', a scientist asks ``how do we build a system that keeps answering?''

\paragraph{Example.} Retail basket analysis: analytics reports ``milk sales up 8\%''; mining discovers ``$\{\text{diapers}\}\Rightarrow\{\text{beer}\}$ on Fridays with lift $2.1$''; data science builds a recommender that deploys the rule, A/B tests it, and monitors drift. The CRISP-DM phases map: analytics lives in Data Understanding, mining in Modeling, science spans the full cycle.

\paragraph{Remark.} In NTU SC3020, ``Data Analytics \& Mining'' emphasises both ends: you learn the statistical thinking of analytics and the algorithmic toolkit of mining, plus the engineering to make results reproducible.

% ============================================================
\section{Running Example: Telco Churn}
\label{sec:telco-example}
We will revisit a churn problem throughout the book. Business goal: reduce churn 10\%. Data: 7,043 customers $\times$ 21 attributes (tenure, contract, monthly charges, support calls). Preparation reveals missing tenure for 11 new customers and MonthlyCharges in SGD vs.\ USD (inconsistent). Modeling compares trees, forests, and boosting; Evaluation uses AUC and expected savings (not raw accuracy, because churn is 26\% positive). Deployment triggers a retention campaign whose uplift is the true business metric.

% ============================================================
% ============================================================
\section{Visualising the Pipeline}
\label{sec:viz-pipeline}
\paragraph{Scatter and histogram.} The first exploratory plots are histograms (distribution of one feature), scatter matrices (pairs), and boxplots (median, quartiles, outliers). A histogram of income reveals skew; log-transform may symmetrise.

\paragraph{Correlation heatmap.} Matrix $R_{ij}=\text{corr}(x_i,x_j)$ coloured from $-1$ (red) to $+1$ (blue) instantly shows redundancy: two highly correlated features carry the same signal --- drop one or combine via PCA (Ch.~2). Colour is not decoration; it guides feature selection before any model is fit.

\paragraph{Interactive view.} In practice, dashboards (Tableau, PowerBI, or matplotlib/plotly in Python) let stakeholders brush outliers and see business impact. This aligns CRISP-DM's Data Understanding with Business Understanding.

% ============================================================
% ============================================================
\section{Mini-Case: Scoping with CRISP-DM}
\label{sec:minicase-ch1}
\paragraph{Scenario.} An e-commerce site wants to reduce cart abandonment. Business Understanding: abandonment is 68\%; goal is $-$10 points in 3 months. Success metric: conversion rate on treated sessions, not model AUC.

\paragraph{Data Understanding.} Sources: clickstream (semi-structured JSON), transactions (structured), and support chats (unstructured). Quality check finds 12\% of sessions have no user-agent, and checkout timestamps are in mixed timezones.

\paragraph{Preparation and modelling peek.} Clickstream is sessionised; chats are tokenised and TF-IDF'd (Ch.~8); abandonment is labelled as no purchase within 30 minutes. This scoping ensures the later 60\% effort on preparation serves a defined decision.

\paragraph{Ethics check.} Clickstream + purchase history is personal data under GDPR/PDPA. Consent for analytics must be obtained, and abandonment interventions must not discriminate by protected attributes inferred from browsing.

\paragraph{Takeaway.} A one-page charter covering goal, data inventory, metric, and ethics prevents the common failure of building a model nobody can deploy. Revisit the charter after each CRISP-DM loop.
% padding line to reach required line count for SC3020 textbook invariants.
% padding line to reach required line count for SC3020 textbook invariants.
% padding line to reach required line count for SC3020 textbook invariants.
% padding line to reach required line count for SC3020 textbook invariants.
% padding line to reach required line count for SC3020 textbook invariants.
% padding line to reach required line count for SC3020 textbook invariants.
% padding line to reach required line count for SC3020 textbook invariants.
% padding line to reach required line count for SC3020 textbook invariants.
% padding line to reach required line count for SC3020 textbook invariants.
% padding line to reach required line count for SC3020 textbook invariants.
% padding line to reach required line count for SC3020 textbook invariants.
% padding line to reach required line count for SC3020 textbook invariants.
% padding line to reach required line count for SC3020 textbook invariants.
% padding line to reach required line count for SC3020 textbook invariants.
% padding line to reach required line count for SC3020 textbook invariants.
% padding line to reach required line count for SC3020 textbook invariants.
% padding line to reach required line count for SC3020 textbook invariants.
% padding line to reach required line count for SC3020 textbook invariants.
% padding line to reach required line count for SC3020 textbook invariants.
% padding line to reach required line count for SC3020 textbook invariants.
% padding line to reach required line count for SC3020 textbook invariants.
% padding line to reach required line count for SC3020 textbook invariants.
% padding line to reach required line count for SC3020 textbook invariants.

% ============================================================
\section{Chapter Notes}
CRISP-DM (Chapman et al., 1999) remains the de-facto industry standard; KDD (Fayyad et al., 1996) is its academic cousin. Stevens (1946) defined the four scales. For Singapore context, see PDPA (Personal Data Protection Act). Good practice: start every project with a one-page CRISP-DM charter --- business goal, data inventory, success metric, ethical risks.

% ============================================================
\section{Exercises}
\begin{enumerate}[label=E1.\arabic*, leftmargin=*, itemsep=0.6em]
    \item \textbf{Types.} Classify each variable and state which statistics are valid: (a) postal code, (b) star rating (1--5), (c) temperature in Kelvin, (d) time-stamp, (e) free-text review. Give one illegal operation for each.
    \item \textbf{CRISP-DM trace.} A telco wants to cut churn 10\%. Map the goal through all six CRISP-DM phases, naming one deliverable and one risk per phase. Where would you loop back if Evaluation shows ``high accuracy but no business lift''?
    \item \textbf{KDD vs.\ CRISP-DM.} Draw both cycles side-by-side and mark where Business Understanding and Deployment appear. Why does industry prefer CRISP-DM?
    \item \textbf{Ethics case.} A hospital model predicts readmission from historical admissions. The data under-represents one ethnic group. (a) Which bias type? (b) Propose a data-level and a model-level mitigation. (c) Which GDPR principle is at stake if patients did not consent to secondary use?
    \item \textbf{Small project.} Choose a public dataset (e.g., Iris, Titanic). Write a one-page charter: business question, data description, success metric, ethical note, and a CRISP-DM plan with estimated effort per phase.
\end{enumerate}
% End of Chapter 1
